% Hand-authored prose; all numbers via macros from generated/macros.tex.

This study crossed five belief-formation rules with three allocators over an
identical feasible set, evaluated all \nTrials{} configurations walk-forward across
\oosPeriods{} out-of-sample weeks net of costs, and tested each against equal
weighting with methods robust to the distributional features that real return
series exhibit.

Three conclusions follow, in decreasing order of confidence.

\begin{enumerate}[noitemsep]
  \item \textbf{Replacing convex optimisation with Gaussian-process search did not
  improve risk-adjusted performance}, changing mean Sharpe by \optimiserDelta{}
  while raising mean turnover from \convexTurnover{} to \gpTurnover{}. Holding the
  objective and the feasible set fixed, the sophisticated optimiser was the more
  expensive way to reach the same place.
  \item \textbf{No configuration separated from the 1/N benchmark}
  (\nBeatingBenchmark{} of \nStrategies{} at the 5\% level). On this sample the
  data cannot distinguish these methods from equal weighting.
  \item \textbf{The choice of estimator moved results less than the choice of
  allocator objective}, though at $N=\nAssets{}$ and $T=\trainWindow{}$ this
  design has limited power to detect the estimation-error reduction that shrinkage
  is built to provide.
\end{enumerate}

The methodological conclusion is more transferable than any of these. The earlier
version of this work reported a large, confident, and entirely spurious effect,
and it did so not through an exotic mistake but through the most ordinary one:
evaluating a model on the data used to fit it. The infrastructure that prevents it
--- a rolling-origin split, a shared constraint object, a test that corrupts the
future and checks the past is unchanged --- amounts to a few hundred lines and
should be written before any method is, not after a result is in hand.

\paragraph{Future work.} In order of expected value:

\begin{enumerate}[noitemsep]
  \item \textbf{Widen the universe.} Fifty to a hundred assets would put $N/T$ in
  the regime where shrinkage estimators are designed to bind, converting the
  weakest limitation of this study into its central experiment.
  \item \textbf{Add a factor-model covariance estimator} (statistical or
  fundamental) as a third point on the estimator axis, since factor structure is
  the standard industrial answer to covariance estimation at scale.
  \item \textbf{Extend to multiple asset classes and regimes}, so the
  out-of-sample window is not a single equity bull market.
  \item \textbf{Model market impact}, under which the turnover difference between
  the two optimisers becomes a first-order rather than second-order effect.
  \item \textbf{Report a full multiple-testing correction across the design}, such
  as the false discovery rate over the whole configuration grid, rather than
  deflation of the winner alone.
\end{enumerate}
